\section{Q\&A 方法论：模型规模到底是什么}

最后的问答把多个 slide 重新串联起来。听众询问视觉模型与语言模型谁更大，Lucas 没有给出单一数字，而是区分参数量、每样本 FLOPs、tokenization 结构与研究资源。这段 teacher voice 很重要，因为它阻止读者把“scale”误写成参数计数排行榜。

\subsection{参数量、FLOPs、内存与数据不是同一尺度}

语言模型的大词表 embedding 会贡献大量参数，视觉 patch projection 的参数结构不同；另一方面，图像 token 数、分辨率和 attention activation 会显著影响 compute。一个模型可以参数较少却 FLOPs 更高，也可以训练 FLOPs 高但推理复用效率好。

\begin{longtable}{L{0.19\textwidth}L{0.31\textwidth}L{0.40\textwidth}}
\toprule
尺度 & 它衡量什么 & 主要盲点 \\
\midrule
参数量 $N$ & 权重存储与部分表示容量 & 不含 activation、token 数、数据和实际利用率 \\
训练 FLOPs & 完成预训练的理论运算量 & 不含低利用率、通信、失败实验和数据成本 \\
每样本 FLOPs & 单次 forward/backward 的算术量 & 不等于 wall-clock latency 或 energy \\
Activation memory & 序列/分辨率、batch 与层数产生的中间状态 & 可被 checkpointing 换成额外计算 \\
数据规模 $D$ & 样本/信息覆盖 & 原始样本数不反映重复、噪声与多样性 \\
下游样本效率 & 表示适配新任务所需标注 & 受 adaptation protocol 与 task similarity 影响 \\
\bottomrule
\end{longtable}

\begin{knowledgebox}{术语消化：capacity、scale 与 efficiency}
\term{capacity} 是模型可表示/拟合函数的能力，不能由参数量完全刻画；\term{scale} 是数据、模型和计算的联合范围；\term{efficiency} 必须指定分母，例如 accuracy per FLOP、per dollar、per joule 或 per labeled example。没有分母的“更高效”没有技术含义。
\end{knowledgebox}

\teachervoice{Lucas 的个人判断是参数量和 FLOPs 都不是“唯一正确”的模型大小。视觉模型当时规模较小，一部分原因是研究兴趣和资源投入差异，而不是视觉任务存在已证明的容量上限。课堂提示：观察到的 frontier 也反映组织与资源分配。}

\subsection{从课堂结果到工程决策}

如果实际任务数据较少、延迟严格，强先验 CNN 或 hybrid 可能更合适；如果拥有大规模多域数据、dense accelerator 与跨任务迁移需求，ViT 类模型可能更有优势；若系统需要高分辨率 dense prediction，还要重新评估 token cost、hierarchy 与 memory。选择模型应从约束反推，而不是从最新榜单顺推。

\begin{importantbox}{工程决策模板}
\begin{enumerate}
  \item 定义 downstream task family、domain shift 与标签预算；
  \item 固定 latency、memory、training compute 与 data governance 约束；
  \item 比较至少一个强 CNN/hybrid 与一个 ViT family baseline；
  \item 同时报告 full-data、few-shot、robustness、calibration 与 failure groups；
  \item 对 pretraining/downstream 做 exact、near 与 semantic dedup audit；
  \item 用真实 hardware profile 验证理论 FLOPs 结论。
\end{enumerate}
\end{importantbox}

\subsection{本章小结}

模型规模没有单一标尺。参数、FLOPs、activation、数据和样本效率回答不同问题；研究 frontier 还受到资源投入影响。工程选择必须从任务与系统约束出发，并保留多种强 baseline。

